\chapter{Splinters}

\section*{Definition}
A splinter is a foreign object that becomes embedded in the skin. Splinters are typically sharp, thin pieces of wood, glass, metal, or plastic. Most are superficial and can be removed at home.

\section*{What Happens in Your Body}
Splinters penetrate the epidermis and dermis, causing a breach in the skin barrier. The body mounts an inflammatory response, which can lead to pain, redness, and swelling. If the splinter is not removed, the body may wall it off with granulation tissue, or secondary bacterial infection (cellulitis) may develop.

\section*{Causes}
\begin{itemize}
\item Wood (most common: from handling lumber, gardening)
\item Metal (wire, staples, metal shavings)
\item Glass (broken glass, fiberglass)
\item Plastic or rubber
\item Thorns or plant spines
\item Fish fins or bones
\item Pencils (graphite can leave a tattoo)
\end{itemize}

\section*{When to See a Doctor}
\begin{itemize}
\item Visible foreign body embedded in the skin
\item Pain at the entry site
\item Redness and mild swelling
\item Sensation of something stuck in the skin
\item Bleeding at the entry point
\item Signs of infection: increased redness, warmth, pus, red streaks
\end{itemize}

\section*{Related Symptoms}
\begin{itemize}
\item Cellulitis (secondary infection)
\item Foreign body granuloma
\item Tetanus risk (especially with soil-contaminated wounds)
\item Puncture wound
\end{itemize}

\section*{Self-Care / Home Management}
\begin{itemize}
\item Wash the area with soap and water
\item Use tweezers sterilized with alcohol to grasp the splinter close to the skin surface
\item Pull the splinter out in the same direction it entered
\item For splinters under the skin, use a sterile needle to gently lift the skin over the splinter
\item After removal, wash again and apply antibiotic ointment
\item Monitor for signs of infection
\end{itemize}

\section*{How Your Doctor Will Evaluate You}
\begin{itemize}
\item Clinical examination and visualization
\item X-ray for radiopaque foreign bodies (metal, glass)
\item Ultrasound for non-radiopaque objects
\end{itemize}

\section*{Treatment Options}
\begin{itemize}
\item Removal with tweezers or sterile needle
\item Local anesthesia for deep or difficult splinters
\item Small incision for deeply embedded splinters
\item Wound irrigation after removal
\item Tetanus preventive treatment if vaccination is not up to date
\item Antibiotics if signs of infection develop
\end{itemize}

\section*{References}
\begin{thebibliography}{99}
\bibitem{splinters:ref1} Capellan O, Hollander JE. ``Management of foreign bodies in the skin.'' \textit{Emergency Medicine Clinics of North America}, 2001;19(3):669--683.
\bibitem{splinters:ref2} Halaas GW. ``Management of foreign bodies in the skin.'' \textit{American Family Physician}, 2007;76(5):683--688.
\bibitem{splinters:ref3} Tully J, Blakeney WG, Adie S. ``Management of foreign body injuries to the hand.'' \textit{Journal of Hand Surgery (European Volume)}, 2017;42(1):87--92.
\bibitem{splinters:ref4} Wronka KS, O'Shea K, Ramesh P, et al. ``Retained foreign bodies in the hand: A review of management.'' \textit{Hand Surgery}, 2004;9(2):145--149.
\bibitem{splinters:ref5} Long B, Koyfman A, Gottlieb M. ``Evaluation and management of superficial foreign bodies in the emergency department.'' \textit{Journal of Emergency Medicine}, 2020;59(4):535--544.
\bibitem{splinters:ref6} Haneke E. ``Splinters and foreign bodies under the nail.'' \textit{Journal of the European Academy of Dermatology and Venereology}, 2010;24(7):754--758.
\end{thebibliography}
